%! Polynomial Equations, Inequalities, and Modeling — Unit 2, Lesson 2.7 — Homework (Answer Key)
\documentclass[10pt]{article}
\usepackage{precalculus-article}
\usepackage{precalculus-key}   % pulls in -boxes; do NOT also load -boxes
\newcommand{\TallMath}[1]{$\displaystyle #1\rule[-1.4em]{0pt}{3.2em}$}

\begin{document}

\pageheader{Unit 2, Lesson 2.7}{Homework --- Answer Key}

\vspace{0.12in}

\begin{enumerate}[itemsep=10pt]

\item Solve each equation by factoring. Set one side to $0$ first.

\begin{enumerate}[label=(\alph*), itemsep=6pt]
  \item $x^3 - 7x^2 + 12x = 0$

        \begin{work}
          x^3 - 7x^2 + 12x &= 0 \\
          x\big(x^2 - 7x + 12\big) &= 0 \\
          x(x-3)(x-4) &= 0
        \end{work}

        Solutions: \ans{$x = 0,\ x = 3,\ x = 4$}

  \item $x^3 = 9x$

        \begin{work}
          x^3 - 9x &= 0 \\
          x\big(x^2 - 9\big) &= 0 \\
          x(x-3)(x+3) &= 0
        \end{work}

        Solutions: \ans{$x = 0,\ x = 3,\ x = -3$}

        Which solution disappears if you divide both sides by $x$ in the first
        step? \quad \ans{$x = 0$}
\end{enumerate}

\item Solve $(x+2)(x-1)(x-5) < 0$ by sign analysis.

\smallskip
\renewcommand{\arraystretch}{1.5}
\begin{center}
\begin{tabular}{|l|c|c|c|c|}
\hline
\textbf{Interval} & $(-\infty,-2)$ & $(-2,1)$ & $(1,5)$ & $(5,\infty)$ \\
\hline
\textbf{Test value} & \ans{$-3$} & \ans{$0$} & \ans{$2$} & \ans{$6$} \\
\hline
\textbf{Sign} & \ans{$-$} & \ans{$+$} & \ans{$-$} & \ans{$+$} \\
\hline
\end{tabular}
\end{center}
\smallskip

Solution in interval notation: \ans{$(-\infty,-2) \cup (1,5)$}

\item Let $p(x) = (x-3)^2(x+1)$. Solve $p(x) \ge 0$.

\begin{enumerate}[label=(\alph*), itemsep=6pt]
  \item Zeros and multiplicities: \ans{$x = 3$ (mult.\ 2), $x = -1$ (mult.\ 1)}
  \item Sign on $(-\infty,-1)$: \ans{$-$} \quad
        on $(-1,3)$: \ans{$+$} \quad
        on $(3,\infty)$: \ans{$+$}
  \item Solution in interval notation: \ans{$[-1, \infty)$}
  \item At which zero does the sign \textit{fail} to change, and what feature
        of that zero explains it? \quad \ans{$x = 3$; even multiplicity}
\end{enumerate}

\boxguard[20]
\item A sheet of card measures 18 cm by 24 cm. Squares of side $x$ are cut from
      the four corners and the flaps folded up to make an open box, so
      $V(x) = x(18-2x)(24-2x)$.

\begin{enumerate}[label=(\alph*), itemsep=6pt]
  \item State the domain restriction the context forces, and say in a few words
        where it comes from. \quad \ans{$0 < x < 9$; the 18 cm side}
  \item Write $V(x)$ in expanded form.

        \begin{work}
          V(x) &= x(18-2x)(24-2x) \\
               &= x\big(432 - 36x - 48x + 4x^2\big) \\
               &= x\big(4x^2 - 84x + 432\big) \\
               &= 4x^3 - 84x^2 + 432x
        \end{work}

        $V(x) = $ \ans{$4x^3 - 84x^2 + 432x$}
  \item Find $V(4)$ and report it in a sentence with units.

        \vspace{0.05in}
        \ansline{$V(4) = 4(10)(16) = 640$: a 4 cm cut makes a box holding $640$ cubic centimeters.}
\end{enumerate}

\item A data set has equally spaced inputs. Its successive differences are the
      route to the right regression.

\smallskip
\renewcommand{\arraystretch}{1.4}
\begin{center}
\begin{tabular}{|c|c|c|c|c|c|c|}
\hline
$x$ & $0$ & $1$ & $2$ & $3$ & $4$ & $5$ \\
\hline
$y$ & $2$ & $2$ & $8$ & $26$ & $62$ & $122$ \\
\hline
\end{tabular}
\end{center}
\smallskip

\begin{enumerate}[label=(\alph*), itemsep=6pt]
  \item First differences: \ans{$0,\ 6,\ 18,\ 36,\ 60$}
  \item Second differences: \ans{$6,\ 12,\ 18,\ 24$}
  \item Third differences: \ans{$6,\ 6,\ 6$}
  \item Which regression --- linear, quadratic, cubic, or quartic --- fits this
        data exactly, and how do the differences tell you?
        \quad \ans{cubic; the third differences are constant}
  \item The calculator returns the coefficients. Name one thing it does
        \textit{not} decide for you. \quad \ans{the degree of the model}
\end{enumerate}

\item The solution set of $x(x-2)^2 > 0$ is

\begin{enumerate}[label=\Alph*., itemsep=3pt]
  \item $(0, \infty)$
  \item $(0,2) \cup (2,\infty)$ \quad \textcolor{keyred}{\textbf{$\leftarrow$ correct: $x=2$ gives $0$, and $0 > 0$ is false}}
  \item $[0, \infty)$
  \item $(2, \infty)$
\end{enumerate}

\end{enumerate}

\vspace{0.1in}

\boxguard
\begin{extensionbox}
Why is one test value per interval enough? Suppose $p$ is positive at some
point of an interval and negative at another point of the same interval, and
that the interval contains no zeros of $p$. Explain, in two or three sentences,
why that is impossible. (Your argument will use the fact that a polynomial
graph has no breaks in it --- calculus will give that fact a name.)

\vspace{0.05in}
\ansline{The graph of a polynomial has no breaks, so to get from a positive output down to a}
\ansline{negative one it has to pass through $0$ somewhere in between --- and that point would}
\ansline{be a zero inside the interval, which we assumed there were none of.}
\end{extensionbox}

\vspace{0.1in}

\boxguard
\begin{notesbox}{Preview: The Unit 2 Test, and Unit 3}
That closes Unit 2. The \textbf{Unit 2 Test} runs the whole arc: end behavior,
real and complex zeros, multiplicity, equivalent forms, and today's equations,
inequalities, and models. Then \textbf{Unit 3, Rational Functions}, divides one
polynomial by another. Today's sign chart shows up again almost immediately
there --- with one new rule. A rational expression can change sign at a zero of
its \textit{denominator} too, at a point that is not even in the domain.
\end{notesbox}

\end{document}
